% Junior Level 2, Week 9 -- workbook chapter (read aloud by a grown-up)
\chapterweek{9}{Project Studio}{Plan on a game card, build the core, then playtest.}

\section{What you will learn}
\begin{ican}
  \item I can plan my game on a game plan card.
  \item I can build my game one must-have at a time.
  \item I can watch a Player test my game and fix a bug.
\end{ican}

\section{My First Arcade Game}
Today you make your own arcade game! A \textbf{must-have} is a part every game needs.

\begin{exercise}[My Game Plan Card]
\begin{enumerate}[leftmargin=1.4em,itemsep=2pt]
  \item Use your rules from last week.
  \item Fill in the game plan card in your project brief.
  \item Draw your stage and your pieces on paper.
  \item A partner plays your paper game. Can they win? Can they lose?
\end{enumerate}
\end{exercise}

\needspace{30\baselineskip}
\begin{exercise}[Build My Game]
Build one must-have. Test it. Check it off. Then the next one!
\begin{center}
\renewcommand{\arraystretch}{1.3}
\begin{tabular}{@{}c l@{}}
\tickbox & A forever loop keeps it going.\\
\tickbox & An if touching makes something happen.\\
\tickbox & A score or misses variable counts.\\
\tickbox & Random gives a surprise.\\
\tickbox & You can win, or it is game over.\\
\end{tabular}
\end{center}
Mei's bananas fall and wait a surprise time:
\begin{center}
\begin{scratch}[scale=0.75]
\blockinit{when \greenflag clicked}
\blockmove{set rotation style \selectmenu{left-right}}
\blockmove{go to \selectmenu{Star}}
\blockmove{point in direction \ovalnum{180}}
\blockinfloop{forever}{
  \blockmove{move \ovalnum{5} steps}
  \blockif{if \boolsensing{touching \ovalsensing*{Bowl} ?} then}{
    \blockvariable{change \selectmenu{score} by \ovalnum{1}}
    \blocksound{start sound \selectmenu{Pop}}
    \blockmove{go to \selectmenu{Star}}
    \blockcontrol{wait \ovaloperator{pick random \ovalnum{1} to \ovalnum{3}} seconds}
  }
  \blockif{if \boolsensing{touching \ovalsensing*{edge} ?} then}{
    \blockvariable{change \selectmenu{misses} by \ovalnum{1}}
    \blockmove{go to \selectmenu{Star}}
    \blockcontrol{wait \ovaloperator{pick random \ovalnum{1} to \ovalnum{3}} seconds}
  }
}
\end{scratch}
\end{center}
\end{exercise}

\needspace{13\baselineskip}
\begin{exercise}[Playtest]
A Player plays your game. Watch. Keep your hands in your lap! Write or draw each bug.
\begin{center}
\renewcommand{\arraystretch}{2}
\begin{tabular}{|p{8.5cm}|c|}
\hline
\textbf{What went wrong?} & \textbf{Fixed?}\\ \hline
 & yes \quad not yet\\ \hline
 & yes \quad not yet\\ \hline
 & yes \quad not yet\\ \hline
\end{tabular}
\end{center}
\end{exercise}

\begin{grownup}[Playtest tip]
The maker says nothing while the Player plays. If the Player is stuck, that is a bug to write down.
\end{grownup}

\begin{vocab}
  \vocabword{prototype}{A quick test version of a game, made before the real one.}
  \vocabword{playtest}{Someone plays your game so you can find bugs.}
\end{vocab}

\begin{reflection}
Color one star for each must-have that works: \quad {\Large\faStar[regular]\ \faStar[regular]\ \faStar[regular]\ \faStar[regular]\ \faStar[regular]}

How do you feel about your game? Circle one.

\begin{center}
\Huge\faSmile[regular]\qquad\faMeh[regular]\qquad\faFrown[regular]
\end{center}
\end{reflection}
